\chapter{Domande di ripasso}

Questo capitolo raccoglie le domande di ripasso proposte durante il corso e una selezione di domande tipiche, organizzate per
argomento, con risposte sintetiche. Vale l'avvertenza data in aula: \textbf{le risposte devono essere complete e possono essere
concise}. Le domande della sessione di ripasso su HTTP (parti di un URL, handshake, pipelining), DNS (cambio di indirizzo,
risoluzione con cache) e CIDR (tavola con L1, L2, R0) sono svolte per esteso nei capitoli di esercizi precedenti.

\info{Tre tipi di domande}{
    \begin{enumerate}
      \item \textbf{Definitorie} (``che cos'è X?''): definizione precisa più almeno un esempio.
      \item \textbf{Di confronto} (``confrontare X e Y''): individuare l'asse lungo cui differiscono (locale/globale, con/senza stato,
            con/senza connessione, hardware/software) e dirlo esplicitamente.
      \item \textbf{Di calcolo}: subnetting, longest prefix matching, CRC, checksum, Dijkstra, frammentazione, RTT, finestre. Vanno
            esercitate a mano: contano i passaggi, non solo il risultato.
    \end{enumerate}
}

% ══════════════════════════════════════════════════════════════
\section{Fondamenti}
% ══════════════════════════════════════════════════════════════

\begin{esercizio}{Reti e Internet}
\begin{enumerate}[label=(\alph*)]
    \item Che cos'è una rete di calcolatori? Che cosa significano ``autonomi'' e ``interconnessi''?
    \item Che cos'è un protocollo?
    \item Confrontare commutazione di circuito e commutazione di pacchetto.
    \item Perché si dice che Internet è una ``rete di reti''? Che ruolo hanno ISP, IXP e reti dei content provider?
    \item Perché IP è la ``vita stretta'' della clessidra di Internet?
\end{enumerate}
\end{esercizio}

\svolgimento
\textbf{(a)} Una collezione di dispositivi di calcolo \textbf{autonomi} (nessuno comanda gli altri: ognuno ha capacità di calcolo
propria) e \textbf{interconnessi} (possono scambiarsi informazione, via cavo, fibra, radio\dots).

\textbf{(b)} Un protocollo definisce il \textbf{formato} e l'\textbf{ordine} dei messaggi scambiati tra due o più entità, e le
\textbf{azioni} da compiere quando un messaggio viene inviato o ricevuto.

\textbf{(c)} Nel \textbf{circuito} le risorse (una banda in FDM, uno slot in TDM) sono \textbf{riservate} per tutta la sessione: prestazioni
garantite, ma capacità sprecata quando l'utente tace e costo di instaurazione. Nel \textbf{pacchetto} i dati sono divisi in pacchetti che
condividono i collegamenti \textbf{su richiesta}, con store-and-forward: nessuna garanzia (code, ritardi, perdite), ma molti più utenti
serviti grazie alla multiplazione statistica (35 utenti invece di 10 nell'esempio classico). È la scelta vincente per il traffico dati, che
è a raffiche.

\textbf{(d)} Perché è l'interconnessione di moltissime reti indipendenti. Gli host si collegano a \textbf{ISP di accesso}, che si collegano
a ISP regionali e questi a pochi ISP globali (tier 1); gli ISP dello stesso livello si scambiano traffico con accordi di
\textbf{peering}, spesso negli \textbf{IXP}, punti neutrali di scambio che evitano di collegare tutti con tutti. I grandi
\textbf{content provider} (Google, Microsoft, \dots) hanno reti proprie che arrivano vicino agli utenti, scavalcando la gerarchia.

\textbf{(e)} Sopra IP ci sono moltissime applicazioni e protocolli di trasporto, sotto moltissime tecnologie di collegamento; in mezzo
c'è \textbf{un solo} protocollo di rete. Tutto ciò che parla IP può comunicare con tutto il resto: nuove applicazioni e nuove tecnologie
fisiche si aggiungono senza cambiare il centro.

\begin{esercizio}{Topologie e categorie di reti}
\begin{enumerate}[label=(\alph*)]
    \item Distinguere collegamenti punto-punto e multipunto.
    \item Confrontare le topologie a bus, anello, stella e maglia (costo, guasti, prestazioni).
    \item Distinguere PAN, LAN, MAN e WAN.
    \item Distinguere comunicazione simplex, half-duplex e full-duplex.
\end{enumerate}
\end{esercizio}

\svolgimento
\textbf{(a)} Punto-punto: un collegamento \textbf{dedicato} tra esattamente due dispositivi. Multipunto (broadcast): più dispositivi
\textbf{condividono} lo stesso collegamento e sentono tutti le trasmissioni altrui, quindi serve un protocollo di accesso al mezzo.

\textbf{(b)}
\begin{itemize}
    \item \textbf{Bus}: economico e semplice, ma la dorsale è un punto di guasto unico e le collisioni aumentano con il traffico.
    \item \textbf{Anello}: ogni nodo ripete il segnale; niente collisioni (con token), ma un guasto può interrompere l'anello (se non è doppio).
    \item \textbf{Stella}: $n$ collegamenti, facile da gestire e da estendere, un guasto isola un solo nodo; il centro è il punto critico.
    \item \textbf{Maglia completa}: massima robustezza e prestazioni, ma $n(n-1)/2$ collegamenti: costo quadratico. Nella pratica si usano
          maglie parziali e topologie ibride.
\end{itemize}

\textbf{(c)} In base all'estensione: \textbf{PAN} (pochi metri, dispositivi di una persona, per esempio Bluetooth); \textbf{LAN}
(edificio o campus, privata, bassa latenza e pochi errori, tempo di trasmissione nel caso peggiore noto); \textbf{MAN} (una città, per
esempio le reti via cavo); \textbf{WAN} (un paese o un continente: host collegati da una subnet di router e linee di trasmissione, spesso
affittate).

\textbf{(d)} Simplex: un solo verso (TV). Half-duplex: entrambi i versi ma a turno (walkie-talkie). Full-duplex: entrambi
contemporaneamente (telefono, Ethernet commutata).

\begin{esercizio}{Il modello a strati}
\begin{enumerate}[label=(\alph*)]
    \item Elencare i 7 strati del modello ISO/OSI e il compito di ciascuno.
    \item Qual è la differenza tra servizio, interfaccia e protocollo?
    \item Distinguere servizi orientati alla connessione e senza connessione.
    \item Spiegare incapsulamento e decapsulamento.
    \item Perché un router non implementa l'intera pila?
    \item Perché la pila TCP/IP si è imposta sul modello OSI?
    \item Quali vantaggi offre l'architettura a strati, e qual è il suo prezzo?
\end{enumerate}
\end{esercizio}

\svolgimento
\textbf{(a)} Dal basso: \textbf{fisico} (trasmissione dei bit sul mezzo), \textbf{collegamento} (frame tra nodi adiacenti, rilevazione
degli errori, accesso al mezzo), \textbf{rete} (instradamento dei pacchetti da host a host), \textbf{trasporto} (da processo a processo,
affidabilità, controllo di flusso), \textbf{sessione} (dialogo e sincronizzazione), \textbf{presentazione} (rappresentazione dei dati:
codifica, compressione, cifratura), \textbf{applicazione} (servizi per l'utente). La pila di Internet ne usa cinque: sessione e
presentazione confluiscono nell'applicazione.

\textbf{(b)} Il \textbf{servizio} è ciò che uno strato offre a quello superiore (\emph{che cosa} fa); l'\textbf{interfaccia} è il modo in cui
lo strato superiore vi accede (le primitive); il \textbf{protocollo} sono le regole con cui lo strato dialoga con la propria entità paritaria
sull'altra macchina (\emph{come} lo fa). Il protocollo si può cambiare senza toccare il servizio.

\textbf{(c)} Orientato alla connessione: si stabilisce una connessione, la si usa, la si chiude (come una telefonata); l'ordine è
preservato. Senza connessione: ogni messaggio porta l'indirizzo completo e viaggia per conto suo (come la posta), l'ordine non è
garantito. Entrambi possono essere affidabili o no.

\textbf{(d)} Scendendo (mittente) ogni strato aggiunge al payload ricevuto il proprio header (ed eventualmente un trailer); risalendo
(destinatario) lo stesso strato lo legge e lo rimuove. Il payload è \textbf{opaco}: uno strato non sa, e non deve sapere, che cosa contiene.

\textbf{(e)} Perché il suo compito è la consegna \textbf{da host a host}: gli bastano i livelli fisico, collegamento e rete. A quale processo
sia destinato il pacchetto non gli interessa; solo gli host implementano tutta la pila.

\textbf{(f)} Per \textbf{tempismo} (TCP/IP era già diffuso quando uscirono i protocolli OSI), \textbf{progetto} (modello e protocolli
complessi, strati squilibriati), \textbf{implementazioni} (le prime erano enormi e lente) e \textbf{politica} (OSI era percepito come
imposto da ministeri e burocrazie), mentre le implementazioni TCP/IP erano gratuite, per esempio in Unix BSD. OSI resta lo standard
\emph{de iure} e un ottimo modello concettuale; TCP/IP è lo standard \emph{de facto}.

\textbf{(g)} Vantaggi: una \textbf{struttura} esplicita con cui ragionare sul sistema, e soprattutto la \textbf{modularità} (si cambia
l'implementazione di uno strato, per esempio un driver o una tecnologia di accesso, senza toccare gli altri, purché il servizio resti
lo stesso). Prezzo: \textbf{maggiore complessità} (interfacce da progettare e implementare, overhead), \textbf{ridondanze} (funzioni come
il controllo degli errori ripetute in più strati), \textbf{rigidità} (le informazioni interne a uno strato sono nascoste; quando servono
altrove bisogna riprogettare o ricorrere a soluzioni \emph{quick and dirty}).

\begin{esercizio}{Banda e throughput}
\begin{enumerate}[label=(\alph*)]
    \item Distinguere transmission rate, banda e throughput.
    \item Che cosa significa store-and-forward? E cut-through?
\end{enumerate}
\end{esercizio}

\svolgimento
\textbf{(a)} Il \textbf{transmission rate} è il massimo che un dispositivo può trasmettere; la \textbf{banda} è il massimo che un
\emph{cammino} (collegamenti e nodi) può trasportare; il \textbf{throughput} è la quantità che sta \emph{effettivamente} fluendo.
Vale sempre throughput $\le$ banda $\le$ rate del dispositivo più lento.

\textbf{(b)} Store-and-forward: un nodo riceve \textbf{tutto} il pacchetto prima di inoltrarlo (può controllarne l'integrità, ma
aggiunge $L/R$ di ritardo a ogni salto). Cut-through: inizia a inoltrare appena ha letto l'intestazione (meno ritardo, ma può propagare
pacchetti danneggiati).

% ══════════════════════════════════════════════════════════════
\section{Livello applicativo}
% ══════════════════════════════════════════════════════════════

\begin{esercizio}{Architetture e protocolli}
\begin{enumerate}[label=(\alph*)]
    \item Confrontare architettura client-server e peer-to-peer.
    \item Perché HTTP è senza stato, e come si recupera lo stato?
    \item Descrivere il caching Web e il GET condizionale.
    \item Perché la posta elettronica usa protocolli diversi?
    \item Descrivere le politiche di BitTorrent.
    \item Perché FTP usa due connessioni? Che differenza c'è tra modalità attiva e passiva?
    \item Qual è la differenza tra un'applicazione e il suo protocollo applicativo?
    \item Quanti indirizzi IP ha un portatile collegato via cavo e via Wi-Fi? A che cosa serve \texttt{127.0.0.1}?
\end{enumerate}
\end{esercizio}

\svolgimento
\textbf{(a)} Nel client-server c'è un \textbf{server} sempre attivo, con indirizzo fisso e noto, e i client non comunicano direttamente tra
loro. Nel P2P i peer sono (idealmente) equivalenti, connessi in modo intermittente, e comunicano direttamente: il sistema scala meglio
perché ogni peer porta anche capacità di servizio, ma è più complesso da gestire e pone problemi di sicurezza e affidabilità.

\textbf{(b)} Per semplicità e per poter servire un numero enorme di richieste senza ricordare nulla dei client. Lo stato si recupera con i
\textbf{cookie}: il server crea un identificativo (\texttt{Set-Cookie}), il browser lo memorizza e lo riallega a ogni richiesta
(\texttt{Cookie}), e il server lo usa come chiave nel proprio database.

\textbf{(c)} Una \textbf{cache} (proxy) soddisfa le richieste al posto del server di origine con copie locali: riduce tempi di risposta e
traffico sul collegamento verso Internet. Con il \textbf{GET condizionale} (\texttt{If-Modified-Since}) la cache chiede l'oggetto solo se
è cambiato; altrimenti il server risponde \texttt{304 Not Modified} senza corpo.

\textbf{(d)} Perché la comunicazione avviene in due tratte diverse: tra mail server si usa \textbf{SMTP} (porta 25, protocollo \emph{push});
per leggere la posta dalla propria casella l'utente usa un protocollo di \emph{accesso} (\textbf{POP3}, \textbf{IMAP} o HTTP per la webmail).

\textbf{(e)} In download: \textbf{rarest first}, si chiedono prima i chunk con meno copie tra i vicini. In upload: \textbf{tit-for-tat}, si
serve chi ci fornisce al ritmo più alto (i migliori 4, rivalutati ogni 10 s), più un vicino scelto a caso ogni 30 s
(\emph{optimistic unchoking}) per dare una possibilità ai nuovi peer.

\textbf{(f)} Per tenere separati comandi e dati (segnalazione \textbf{fuori banda}): la connessione di \textbf{controllo} (porta 21)
resta aperta per tutta la sessione, mentre ogni trasferimento (\texttt{get}, \texttt{put}, anche \texttt{ls}) apre e chiude una
connessione \textbf{dati}. In modalità \textbf{attiva} è il server ad aprire la connessione dati verso il client, cosa impossibile se il
client è dietro un NAT o un firewall; in modalità \textbf{passiva} il client apre entrambe le connessioni.

\textbf{(g)} L'applicazione comprende tutto: per il Web, il formato dei documenti (HTML), i browser, i server. Il protocollo (HTTP)
fissa solo il formato e l'ordine dei messaggi tra le parti, ed è l'unica cosa su cui i due capi devono concordare.

\textbf{(h)} Due: un indirizzo appartiene a un'\textbf{interfaccia}, non alla macchina, e per il livello di rete il portatile è come
due host distinti raggiunti per cammini diversi. \texttt{127.0.0.1} (\texttt{localhost}) indica la macchina stessa: due processi
locali comunicano attraverso di esso come in rete, senza che nulla esca dall'host; serve per provare client e server su un solo computer.

\begin{esercizio}{Socket e DNS}
\begin{enumerate}[label=(\alph*)]
    \item Quali sono le differenze nella sequenza di chiamate tra un'applicazione UDP e una TCP?
    \item Da quanti elementi è identificato un socket UDP? E uno TCP?
    \item Perché il server TCP ha bisogno di due socket?
    \item Confrontare query ricorsive e iterative nel DNS.
\end{enumerate}
\end{esercizio}

\svolgimento
\textbf{(a)}
\begin{center}
\small
\begin{tabular}{p{6.5cm}p{6.5cm}}
\toprule
\textbf{UDP} & \textbf{TCP} \\
\midrule
\texttt{socket()}; \texttt{bind()} solo sul server & \texttt{socket()}; server: \texttt{bind()}, \texttt{listen()}, \texttt{accept()} \\
nessuna connessione & client: \texttt{connect()} (three-way handshake) \\
\texttt{sendto()} / \texttt{recvfrom()} con l'indirizzo a ogni messaggio & \texttt{send()} / \texttt{read()}: l'indirizzo è fissato dalla connessione \\
\texttt{close()} & \texttt{close()} (sul server anche del socket dedicato) \\
\bottomrule
\end{tabular}
\end{center}

\textbf{(b)} Un socket UDP è identificato dalla coppia (IP, porta) di destinazione: lo stesso socket riceve da client diversi. Un socket TCP
connesso dalla quadrupla (IP sorgente, porta sorgente, IP destinazione, porta destinazione).

\textbf{(c)} Il \textbf{welcoming socket} resta in ascolto per accettare nuove connessioni; per ogni client \texttt{accept()} crea un
\textbf{socket dedicato}, usato per la comunicazione. Così il server continua ad accettare client mentre serve quelli già connessi.

\textbf{(d)} In una query \textbf{ricorsiva} il server interrogato si fa carico dell'intera risoluzione e restituisce la risposta finale; in
una \textbf{iterativa} risponde con ciò che sa (spesso un rinvio al server successivo). Tipicamente l'host fa una query ricorsiva al server
locale, che procede in modo iterativo; i server in alto nella gerarchia non offrono ricorsione, per non tenere stato.

% ══════════════════════════════════════════════════════════════
\section{Livello di trasporto}
% ══════════════════════════════════════════════════════════════

\begin{esercizio}{UDP e TCP}
\begin{enumerate}[label=(\alph*)]
    \item Quali sono le responsabilità del livello di trasporto? Quali copre UDP?
    \item Perché usare UDP invece di TCP?
    \item Qual è la relazione tra MSS e MTU?
    \item Perché il fast retransmit aspetta 3 ACK duplicati e non uno?
\end{enumerate}
\end{esercizio}

\svolgimento
\textbf{(a)} (1) Consegna da processo a processo (multiplexing e demultiplexing con le porte); (2) controllo di integrità (checksum);
(3) trasferimento affidabile; (4) controllo di flusso e di congestione. UDP copre solo le prime due; TCP tutte.

\textbf{(b)} Controllo più fine da parte dell'applicazione (utile nel tempo reale), nessun handshake (una query DNS è una domanda e una
risposta), nessuno stato di connessione (più client serviti), header minimo (8 byte contro 20).

\textbf{(c)} $\text{MSS} = \text{MTU} - 40$ byte (20 di header IP e 20 di header TCP, senza opzioni): con MTU Ethernet di 1500 byte,
MSS di 1460 byte.

\textbf{(d)} Perché un ACK duplicato può nascere anche da un semplice \textbf{riordino}: se due segmenti arrivano scambiati, il ricevente
invia un duplicato prima di ricevere quello mancante. Tre duplicati sono un indizio molto più affidabile di una perdita.

\begin{esercizio}{Trasferimento affidabile}
\begin{enumerate}[label=(\alph*)]
    \item Confrontare Go-Back-N e Selective Repeat.
    \item Perché in Selective Repeat la finestra non può superare metà dello spazio dei numeri di sequenza?
    \item Perché il timeout deve essere maggiore dell'RTT, ma non troppo?
\end{enumerate}
\end{esercizio}

\svolgimento
\textbf{(a)} In \textbf{GBN} il ricevente accetta solo pacchetti in ordine e invia ACK cumulativi; il mittente, allo scadere del timer
(unico, sul pacchetto più vecchio), ritrasmette tutta la finestra. Ricevente semplicissimo, ma molte ritrasmissioni inutili. In
\textbf{SR} il ricevente bufferizza i pacchetti fuori ordine e conferma ciascuno; il mittente ha un timer per pacchetto e ritrasmette solo
quelli persi. Più efficiente, più complesso.

\textbf{(b)} Perché altrimenti, se gli ACK vanno persi, un pacchetto ritrasmesso può cadere nella nuova finestra del ricevente ed essere
scambiato per un pacchetto nuovo con lo stesso numero (si veda l'esercizio nel capitolo sul trasporto).

\textbf{(c)} Troppo corto: ritrasmissioni inutili, che aggiungono traffico proprio quando la rete è lenta. Troppo lungo: si reagisce tardi
alle perdite. $\text{EstimatedRTT} + 4\,\text{DevRTT}$ dà un margine che si adatta alla variabilità.

\begin{esercizio}{Connessione e congestione}
\begin{enumerate}[label=(\alph*)]
    \item Descrivere il problema dei due eserciti e la sua rilevanza per TCP.
    \item Distinguere controllo di flusso e controllo di congestione.
    \item Perché un timeout riporta cwnd a 1 MSS mentre 3 ACK duplicati la dimezzano soltanto?
    \item Perché i numeri di sequenza iniziali sono casuali?
\end{enumerate}
\end{esercizio}

\svolgimento
\textbf{(a)} Due eserciti alleati devono attaccare insieme, ma comunicano con messaggeri che possono essere catturati: qualunque sia il
numero di conferme, chi invia l'ultima non sa se è arrivata. Nessun protocollo garantisce l'accordo su un canale inaffidabile. TCP accetta
un compromesso: il three-way handshake (e la chiusura con TIME\_WAIT) non è perfetto ma è adeguato, e i timer risolvono le situazioni
ambigue.

\textbf{(b)} Il \textbf{flusso} protegge il \textbf{ricevente} dal traboccamento del buffer (rwnd, annunciata dal ricevente); la
\textbf{congestione} protegge la \textbf{rete} (cwnd, stimata dal mittente osservando perdite e ritardi). Il mittente rispetta
$\text{LastByteSent} - \text{LastByteAcked} \le \min\{\text{cwnd}, \text{rwnd}\}$.

\textbf{(c)} Un timeout significa che non arriva più nulla: congestione grave, si riparte da capo. Tre ACK duplicati dimostrano che i
segmenti successivi stanno arrivando: congestione lieve, basta dimezzare. In entrambi i casi ssthresh $=$ cwnd/2.

\textbf{(d)} Per non confondere segmenti di una vecchia connessione (con le stesse porte) con quelli di una nuova, e per rendere difficile a
un attaccante indovinare i numeri e iniettare segmenti falsificati.

% ══════════════════════════════════════════════════════════════
\section{Livello di rete}
% ══════════════════════════════════════════════════════════════

\begin{esercizio}{IP e indirizzamento}
\begin{enumerate}[label=(\alph*)]
    \item Che cos'è il servizio best effort e perché è sufficiente?
    \item Distinguere forwarding e routing, data plane e control plane.
    \item Perché un indirizzo IP appartiene a un'interfaccia e non a un host?
    \item Perché l'indirizzamento a classi è stato abbandonato?
    \item Perché IPv6 ha eliminato il checksum e la frammentazione nei router?
\end{enumerate}
\end{esercizio}

\svolgimento
\textbf{(a)} La rete ``fa del suo meglio'' senza garanzie su consegna, ordine, ritardo o banda. Basta perché le garanzie che servono vengono
ricostruite agli estremi (TCP, applicazioni): è il principio \textbf{end-to-end}, nucleo semplice e intelligenza ai bordi.

\textbf{(b)} Il \textbf{forwarding} (data plane) è l'azione locale di spostare un pacchetto dall'ingresso all'uscita giusta: nanosecondi, in
hardware. Il \textbf{routing} (control plane) calcola i cammini sull'intera rete e riempie le tabelle: secondi, in software.

\textbf{(c)} Perché ogni interfaccia è collegata a una subnet diversa e deve inviare e ricevere datagrammi su quella subnet: un router con tre
interfacce ha tre indirizzi, uno per subnet.

\textbf{(d)} Per lo spreco: le classi avevano dimensioni fisse (A 16 milioni, B 65\,536, C 256 indirizzi). A chi servivano 300 indirizzi
toccava una classe B, sprecandone oltre 65\,000. Con CIDR basta un /23 (512 indirizzi).

\textbf{(e)} Il checksum va ricalcolato a ogni salto (il TTL cambia) ed è ridondante, perché collegamento e trasporto controllano già
gli errori. La frammentazione nei router è lenta e complessa: in IPv6 un router che riceve un datagramma troppo grande lo scarta e invia
un ICMPv6 \emph{Packet Too Big}, e sarà il mittente a ridurre la dimensione.

\begin{esercizio}{DHCP, NAT e ICMP}
\begin{enumerate}[label=(\alph*)]
    \item Descrivere i 4 passi di DHCP e spiegare perché usa il broadcast.
    \item Perché il NAT cambia la porta sorgente? Qual è il suo limite fondamentale?
    \item Come funziona traceroute?
\end{enumerate}
\end{esercizio}

\svolgimento
\textbf{(a)} Discover (broadcast, da 0.0.0.0 porta 68 a 255.255.255.255 porta 67), Offer (con l'indirizzo proposto in \texttt{yiaddr}),
Request (il client accetta un'offerta), ACK (conferma). Il broadcast serve perché l'host non ha ancora un indirizzo e non sa chi sia il
server; la request è in broadcast per informare tutti i server di quale offerta è stata scelta.

\textbf{(b)} Perché host diversi possono scegliere la stessa porta: il NAT deve garantire coppie (IP WAN, porta) uniche per distinguere le
risposte. Il limite è che le connessioni \textbf{iniziate dall'esterno} non trovano una riga nella tabella: servono port forwarding, UPnP,
STUN/TURN o hole punching. Inoltre il NAT viola la stratificazione (un dispositivo di rete modifica le porte, che sono di trasporto).

\textbf{(c)} Invia sonde (UDP verso porte improbabili) con TTL $= 1, 2, 3, \dots$ Il router a cui il TTL scade scarta la sonda e risponde con
ICMP \emph{Time Exceeded}, rivelando il proprio indirizzo; quando la sonda raggiunge la destinazione, questa risponde con ICMP \emph{Port
Unreachable} e traceroute si ferma. Per ogni salto misura il tempo di andata e ritorno; un asterisco indica una risposta mancata.

\begin{esercizio}{Routing}
\begin{enumerate}[label=(\alph*)]
    \item Confrontare Distance Vector e Link State.
    \item Spiegare il count-to-infinity e le sue cause.
    \item Perché il flooding, pur trovando sempre il cammino, è impraticabile?
\end{enumerate}
\end{esercizio}

\svolgimento
\textbf{(a)}
\begin{center}
\small
\begin{tabular}{lll}
\toprule
& \textbf{Distance Vector} & \textbf{Link State} \\
\midrule
Conoscenza & locale (vettori dei vicini) & globale (tutto il grafo) \\
Algoritmo & Bellman-Ford & Dijkstra \\
Funzionamento & decentralizzato, asincrono & ogni nodo calcola sulla mappa completa \\
Messaggi & solo ai vicini & link-state packet in flooding a tutti \\
Convergenza & lenta, count-to-infinity & più rapida \\
Robustezza & un errore si propaga ai calcoli degli altri & ogni nodo calcola la propria tabella \\
\bottomrule
\end{tabular}
\end{center}

\textbf{(b)} Il vettore $D_x(y)$ si calcola dai vettori dei vicini, che contengono solo i \textbf{costi} e non i \textbf{cammini}: un nodo
può scegliere un vicino il cui cammino migliore passa per il nodo stesso. Quando un collegamento peggiora, i due nodi si rimandano stime
crescenti a piccoli passi finché non si raggiunge il costo alternativo (o all'infinito, se il collegamento cade). Contromisure: poisoned
reverse, split horizon, un ``infinito'' finito (16 in RIP).

\textbf{(c)} Esplora tutti i cammini ed è robustissimo, ma genera una quantità di traffico enorme (che cresce esponenzialmente senza
controlli sui duplicati): va bene solo in contesti particolari (diffusione di informazioni a tutti, reti militari), non come routing di Internet.

% ══════════════════════════════════════════════════════════════
\section{Livello di collegamento}
% ══════════════════════════════════════════════════════════════

\begin{esercizio}{Framing ed errori}
\begin{enumerate}[label=(\alph*)]
    \item Elencare i metodi di framing e spiegare perché il conteggio dei byte non si usa.
    \item Perché la parità singola è quasi inutile nelle reti reali?
    \item Che cosa rileva e corregge la parità bidimensionale?
    \item Quanti errori rileva e corregge un codice con distanza di Hamming $d$?
\end{enumerate}
\end{esercizio}

\svolgimento
\textbf{(a)} Conteggio dei byte, byte di flag con byte stuffing, bit di flag con bit stuffing, violazioni della codifica fisica. Il
conteggio non si usa (da solo) perché se il campo di conteggio si corrompe non è più possibile ritrovare l'inizio dei frame successivi.

\textbf{(b)} Rileva solo un numero dispari di errori, e gli errori reali arrivano a raffiche: la probabilità di non accorgersene è vicina al 50\%.

\textbf{(c)} Rileva e \textbf{corregge} un errore singolo (incrocio di riga e colonna); rileva ma non corregge qualsiasi coppia di errori.

\textbf{(d)} Rileva $d - 1$ errori e ne corregge $\lfloor (d-1)/2 \rfloor$: servono distanza $k + 1$ per rilevarne $k$, $2k + 1$ per correggerne $k$.

\begin{esercizio}{Accesso multiplo}
\begin{enumerate}[label=(\alph*)]
    \item Quali sono i requisiti ideali di un protocollo ad accesso multiplo?
    \item Qual è il difetto strutturale di TDMA e FDMA?
    \item Se tutti i nodi ascoltano il canale prima di trasmettere, perché ci sono ancora collisioni?
    \item Perché Ethernet ha una dimensione minima dei frame?
    \item Che cosa sono i problemi del terminale nascosto ed esposto, e come li affronta MACA?
\end{enumerate}
\end{esercizio}

\svolgimento
\textbf{(a)} Con un solo nodo attivo, usare tutto il rate $R$; con $M$ nodi attivi, $R/M$ ciascuno in media; essere
\textbf{decentralizzato} (niente punto di guasto unico, niente sincronizzazione globale); essere \textbf{semplice} ed economico.

\textbf{(b)} Un nodo ha sempre e solo $R/N$, anche quando tutti gli altri tacciono: viola il primo requisito.

\textbf{(c)} Per il \textbf{ritardo di propagazione}: un nodo può trovare libero un canale su cui un altro ha già iniziato a trasmettere, ma il
cui segnale non è ancora arrivato.

\textbf{(d)} Perché CSMA/CD rilevi la collisione, il mittente deve essere ancora in trasmissione quando il segnale della collisione torna
indietro: il frame deve durare almeno $2\tau$. Da qui i 64 byte minimi e il campo pad.

\textbf{(e)} Nel wireless \textbf{A} e \textbf{C} possono non sentirsi pur raggiungendo entrambi \textbf{B}: ascoltando il canale lo
trovano libero e collidono in B (\textbf{terminale nascosto}); viceversa un nodo può rinunciare a trasmettere perché sente una trasmissione
che in realtà non disturberebbe il suo destinatario (\textbf{terminale esposto}). MACA fa precedere i dati da un breve \textbf{RTS}, a cui il
destinatario risponde con un \textbf{CTS} che indica la durata: chi sente il CTS tace, chi sente solo l'RTS può trasmettere.

\begin{esercizio}{Switch e ARP}
\begin{enumerate}[label=(\alph*)]
    \item Confrontare switch e router.
    \item Che cos'è il self-learning?
    \item Perché le LAN commutate devono avere una topologia ad albero (senza cicli)?
    \item Come cambiano IP e MAC di destinazione lungo un percorso con più salti?
    \item Perché servono sia gli indirizzi MAC sia gli indirizzi IP?
\end{enumerate}
\end{esercizio}

\svolgimento
\textbf{(a)} Lo \textbf{switch} lavora al livello 2 con indirizzi MAC, è plug-and-play, veloce e senza algoritmo di routing (tabelle
per self-learning). Il \textbf{router} lavora al livello 3 con indirizzi IP, esegue algoritmi di routing, va configurato, isola i domini di
broadcast e supporta qualunque topologia.

\textbf{(b)} Quando riceve un frame, lo switch registra che il \textbf{MAC sorgente} è raggiungibile dalla porta di arrivo; per la
destinazione inoltra se la conosce, fa flooding se non la conosce, filtra se è sulla stessa porta. Le righe scadono dopo un TTL.

\textbf{(c)} Senza routing, un frame in broadcast (o in flooding) in un ciclo girerebbe per sempre moltiplicandosi: una \textbf{tempesta di
broadcast}. Per questo le reti commutate usano alberi (in pratica calcolati da un protocollo di spanning tree).

\textbf{(d)} Il MAC di destinazione è quello del \textbf{prossimo salto} e cambia a ogni collegamento; l'IP di destinazione è quello della
\textbf{destinazione finale} e resta invariato (salvo NAT). Ogni router decapsula e reincapsula il datagramma in un nuovo frame.

\textbf{(e)} Il MAC è piatto, assegnato dal costruttore e portabile: identifica l'interfaccia su un collegamento. L'IP è gerarchico e
dipende dalla rete a cui si è collegati: permette di aggregare e instradare. Tenerli separati rende gli strati indipendenti (una LAN
può trasportare anche protocolli diversi da IP) e consente a una scheda di cambiare rete senza cambiare MAC.

% ══════════════════════════════════════════════════════════════
\section{Sicurezza}
% ══════════════════════════════════════════════════════════════

\begin{esercizio}{Attacchi e crittografia}
\begin{enumerate}[label=(\alph*)]
    \item Distinguere virus e worm.
    \item Elencare i tre tipi di attacco DoS.
    \item Perché lo sniffing è così difficile da rilevare?
    \item Enunciare il principio di Kerckhoffs e spiegarne la ragione.
    \item Perché il one-time pad, pur essendo inviolabile, non si usa in pratica?
    \item Confrontare crittografia simmetrica e asimmetrica.
\end{enumerate}
\end{esercizio}

\svolgimento
\textbf{(a)} Un \textbf{virus} richiede un'interazione dell'utente (aprire un allegato, eseguire un programma); un \textbf{worm} entra da solo,
sfruttando applicazioni di rete vulnerabili, e si propaga automaticamente.

\textbf{(b)} Attacco a una \textbf{vulnerabilità} (messaggi costruiti per far bloccare un servizio o il sistema operativo);
\textbf{bandwidth flooding} (inondare la vittima di pacchetti); \textbf{connection flooding} (aprire moltissime connessioni TCP, anche
semi-aperte come nel SYN flood).

\textbf{(c)} Perché è \textbf{passivo}: non inietta traffico, quindi dall'esterno non lascia tracce. La difesa non è la rilevazione ma la
\textbf{cifratura}: se non si può impedire l'ascolto, si rende incomprensibile ciò che viene ascoltato.

\textbf{(d)} Gli algoritmi devono essere \textbf{pubblici}; solo le \textbf{chiavi} sono segrete. Un algoritmo pubblico viene analizzato da
tutta la comunità e le sue debolezze emergono e si correggono; un algoritmo segreto sembra sicuro solo finché qualcuno non lo rompe, e non
lo si può sapere. La \emph{security through obscurity} non funziona.

\textbf{(e)} Il pad deve essere casuale, lungo quanto i messaggi e mai riusato: va distribuito in modo sicuro, e se esiste un canale sicuro
abbastanza capiente per il pad tanto vale usarlo per il messaggio. In più ci sono problemi di sincronizzazione.

\textbf{(f)} Nella \textbf{simmetrica} (DES, AES) mittente e destinatario condividono la stessa chiave: veloce, ma la chiave va scambiata in
modo sicuro e ne serve una per ogni coppia. Nell'\textbf{asimmetrica} (RSA) ognuno ha una chiave pubblica e una privata: niente segreto da
scambiare, ma molto più lenta. In pratica si usano insieme: l'asimmetrica per scambiare una \textbf{chiave di sessione} simmetrica.

\begin{esercizio}{Autenticazione e SSL}
\begin{enumerate}[label=(\alph*)]
    \item Perché cifrare la password non basta ad autenticarsi?
    \item A che cosa serve un'autorità di certificazione?
    \item Descrivere le fasi di SSL/TLS.
    \item Perché SSL mette un numero di sequenza dentro la parte protetta dei record?
\end{enumerate}
\end{esercizio}

\svolgimento
\textbf{(a)} Perché la password cifrata può essere registrata e \textbf{rigiocata} (playback attack). Serve un \textbf{nonce}: un valore
usato una sola volta, scelto dal verificatore, da cifrare con la chiave condivisa.

\textbf{(b)} Garantisce che una chiave pubblica appartenga davvero a chi dice di esserne il proprietario: firma un \textbf{certificato} che lega
l'identità alla chiave. Senza, un man-in-the-middle potrebbe sostituire la propria chiave (come nell'attacco a Diffie-Hellman).

\textbf{(c)} \textbf{Handshake}: dopo la connessione TCP, il server invia il certificato; il client lo verifica, genera un \emph{master secret},
lo cifra con la chiave pubblica del server e lo invia, e i due si scambiano MAC di tutti i messaggi dell'handshake (contro manomissioni).
\textbf{Derivazione delle chiavi}: dal master secret si ricavano quattro chiavi (cifratura e MAC, per ciascun verso). \textbf{Trasferimento
dati}: il flusso è diviso in record; a ogni record si aggiunge un MAC, si cifra e si passa a TCP.

\textbf{(d)} Il numero di sequenza TCP è in chiaro e un man-in-the-middle potrebbe riordinare, duplicare o eliminare record. Includendo un
numero di sequenza nel calcolo del MAC, ogni alterazione dell'ordine invalida il MAC: si protegge l'integrità dell'intero flusso.

\begin{esercizio}{Firewall e IDS}
\begin{enumerate}[label=(\alph*)]
    \item Qual è il limite del filtraggio tradizionale e come lo risolve quello stateful?
    \item Quando serve un application gateway?
    \item Distinguere IDS basati su firme e basati su anomalie.
    \item A che cosa serve una DMZ?
\end{enumerate}
\end{esercizio}

\svolgimento
\textbf{(a)} Il filtro tradizionale è \textbf{senza stato}: una regola che lascia entrare i segmenti con ACK $= 1$ e porta sorgente 80
fa passare anche segmenti che non appartengono ad alcuna connessione (usati negli attacchi). Il filtro \textbf{stateful} tiene una tabella
delle connessioni (riconosce SYN, SYN-ACK, FIN, e usa timeout di inattività) e ammette solo i pacchetti di connessioni esistenti.

\textbf{(b)} Quando le decisioni dipendono da \textbf{utenti} o \textbf{contenuti applicativi} (per esempio Telnet solo per alcuni utenti
autenticati): queste informazioni non sono negli header IP/TCP/UDP. Il prezzo: un gateway per applicazione, prestazioni più basse, client
configurati per usarlo.

\textbf{(c)} Basati su \textbf{firme}: confrontano il traffico con un database di attacchi noti; precisi, ma ciechi agli attacchi nuovi. Basati
su \textbf{anomalie}: costruiscono un profilo del traffico normale e segnalano ciò che se ne discosta; possono scoprire attacchi nuovi, ma
producono falsi positivi.

\textbf{(d)} È una zona a sicurezza ridotta, tra un firewall esterno più permissivo e uno interno restrittivo, dove si mettono i server che
devono essere raggiungibili da Internet (Web, posta, DNS). Se uno di essi viene compromesso, l'attaccante è ancora fuori dalla rete interna.

% ══════════════════════════════════════════════════════════════
\section{Consigli per la preparazione}
% ══════════════════════════════════════════════════════════════

\warning{Gli errori più frequenti}{
    \begin{itemize}
      \item Confondere consegna \textbf{host-to-host} (rete) e \textbf{process-to-process} (trasporto).
      \item Attribuire il three-way handshake a HTTP invece che a TCP.
      \item Confondere il \textbf{MAC address} (livello 2) con il \textbf{MAC crittografico} (\emph{Message Authentication Code}).
      \item Dimenticare che il MAC di destinazione cambia a ogni salto, l'IP di destinazione no.
      \item Dire che il TTL di IP è un tempo: è un \textbf{contatore di salti}.
      \item Nel CRC: dimenticare gli $r$ zeri, allineare male $G$, usare la sottrazione con prestiti invece dello XOR.
      \item Nel checksum: dimenticare il wrap around o il complemento finale.
      \item Nel subnetting: proporre maschere con bit non contigui; dimenticare che rete e broadcast non sono assegnabili.
      \item Nella frammentazione: esprimere l'offset in byte invece che in unità da 8 byte, o dimenticare che i dati devono essere multipli di 8.
      \item Nella stima dell'RTT: non dichiarare l'ordine di aggiornamento di EstimatedRTT e DevRTT.
    \end{itemize}
}

\info{Come prepararsi}{
    \begin{itemize}
      \item Per ogni capitolo, rileggere il riquadro ``In breve'' e provare a spiegarlo ad alta voce senza guardare.
      \item Rifare gli esercizi di calcolo \textbf{a mano}, su carta, controllando i passaggi con le soluzioni solo alla fine.
      \item Ripetere almeno una volta le esercitazioni con Wireshark: vedere i pacchetti veri fissa i concetti molto più della teoria.
      \item Per le domande di confronto, preparare piccole tabelle a due colonne (come quelle di questo capitolo).
    \end{itemize}
}
